\documentclass[]{article}

\usepackage{amsmath}
\usepackage{multirow}
\usepackage{natbib}
\usepackage{graphicx} 


\begin{document}

The kinetic Sunyaev-Zel'dovich (kSZ) effect provides a unique probe of the baryonic content in galaxy clusters through the scaling relation between Thomson optical depth ($\tau$) and halo mass ($M$), but its faint signal is obscured by dominant Cosmic Microwave Background (CMB) anisotropies and instrumental noise. We test a methodology to constrain this $\tau-M$ relation using a simulated 100 deg$^2$ CMB map, characteristic of current surveys with a 1.4' beam and 20 $\mu$K white noise, and an associated catalog of 5,000 massive halos. Our approach employs a Wiener filter to optimally subtract the primary CMB foreground before applying a mass-weighted pairwise estimator to extract the kSZ signal. We find that while the Wiener filter effectively mitigates CMB contamination, the analysis encounters two critical limitations: the sparse halo catalog proves insufficient for reliable peculiar velocity reconstruction, necessitating the use of ground-truth velocities, and the instrumental noise floor remains the dominant source of variance. Consequently, we report a marginal detection of the kSZ signal at 1.56$\sigma$ significance, which leads to a weak constraint on the scaling relation, yielding a power-law slope of $0.38 \pm 7.23$. While this result is statistically consistent with the theoretical expectation of $2/3$, the large uncertainty demonstrates that for the given survey parameters, constraining the baryonic physics of galaxy clusters is fundamentally limited by instrumental noise and the availability of dense, overlapping catalogs for velocity reconstruction.

\end{document}
                